\SourceLeaf{35}
Mit diesen Untersuchungen ist das Ziel des Abschnitts erreicht. Wir fügen noch einige Sätze hinzu, die auf ähnlichen Grundsätzen beruhen und im Folgenden oft gebraucht werden.

\Needspace{6\baselineskip}
\begin{center}\emph{Verschiedene Sätze.}\end{center}

\Article{38}
\textsc{Aufgabe.} \emph{Man bestimme, wie viele positive Zahlen kleiner als eine gegebene positive Zahl $A$ und zugleich zu ihr teilerfremd sind.}

Der Kürze halber bezeichnen wir die Anzahl der positiven Zahlen, die zu einer gegebenen Zahl teilerfremd und kleiner als sie sind, durch das vorangestellte Zeichen $\phi$. Gesucht ist also $\phi A$.

I. Ist $A$ eine Primzahl, so sind offenbar alle Zahlen von $1$ bis $A-1$ zu $A$ teilerfremd. In diesem Fall ist daher
\[\phi A=A-1\]

II. Ist $A$ eine Primzahlpotenz, etwa $=p^m$, so sind alle durch $p$ teilbaren Zahlen nicht zu $A$ teilerfremd, die übrigen dagegen schon. Von den $p^m-1$ Zahlen müssen also die Zahlen $p,2p,3p....(p^{m-1}-1)p$ wegfallen. Es bleiben $p^m-1-(p^{m-1}-1)$ oder $p^{m-1}(p-1)$. Somit ist
\[\phi p^m=p^{m-1}(p-1)\]

III. Die übrigen Fälle lassen sich mithilfe des folgenden Satzes leicht auf diese zurückführen: \emph{Ist $A$ in paarweise teilerfremde Faktoren $M,N,P$ usw. zerlegt, so ist}
\[\phi A=\phi M.\,\phi N.\,\phi P\ \textit{usw.}\]
Der Beweis lautet so: Die zu $M$ teilerfremden und unter $M$ liegenden Zahlen seien $m,m'.m''$ usw.; ihre Anzahl ist also $=\phi M$. Ebenso seien die zu $N,P$ usw. jeweils teilerfremden und kleineren Zahlen $n,n',n''$ usw., $p,p',p''$ usw.; ihre Anzahlen sind $\phi N,\phi P$ usw. Nun ist jede zum Produkt $A$ teilerfremde Zahl auch zu jedem einzelnen Faktor $M,N,P$ usw. teilerfremd, und umgekehrt (Artikel 19). Ferner ist jede Zahl, die zu einer der Zahlen $m,m',m''$ usw. nach dem Modul $M$ kongruent ist, zu $M$ teilerfremd, und umgekehrt; ebenso gilt dies für $N,P$ usw. Die Frage ist somit darauf zurückgeführt, wie viele Zahlen unter $A$ es gibt, die nach dem Modul $M$ zu einer der Zahlen $m,m',m''$ usw., nach $N$ zu einer der Zahlen $n,n',n''$
\SourceLeaf{36}
usw. kongruent sind. Nach Artikel 32 sind alle Zahlen, die nach den einzelnen Moduln $M,N,P$ usw. bestimmte Reste ergeben, nach deren Produkt $A$ kongruent. Unter $A$ gibt es also nur eine Zahl, die nach den einzelnen Moduln $M,N,P$ usw. zu gegebenen Resten kongruent ist. Die gesuchte Anzahl ist daher gleich der Anzahl der Kombinationen jeder der Zahlen $m,m',m''$ mit jeder der Zahlen $n,n',n''$ und jeder der Zahlen $p,p',p''$ usw. Diese Anzahl ist nach der Kombinationstheorie gleich $=\phi M.\,\phi N.\,\phi P$ usw. \mbox{\emph{Q. E. D.}}

IV. Wie dies auf den vorliegenden Fall anzuwenden ist, sieht man nun leicht. Man zerlege $A$ in seine Primfaktoren oder bringe es in die Form $a^\alpha b^\beta c^\gamma$ usw., wobei $a,b,c$ usw. verschiedene Primzahlen bezeichnen. Dann ist
\[
\begin{aligned}
\phi A&=\phi a^\alpha.\,\phi b^\beta.\,\phi c^\gamma\ \mathrm{usw.}\\
&=a^{\alpha-1}(a-1)\,b^{\beta-1}(b-1)\,c^{\gamma-1}(c-1)\ \mathrm{usw.}
\end{aligned}
\]
oder übersichtlicher
\[\phi A=A\,\frac{a-1}{a}.\,\frac{b-1}{b}.\,\frac{c-1}{c}\ \mathrm{usw.}\]

\emph{Beispiel.} Sei $A=60=2^2.\,3.\,5$, also $\phi A=\frac{1}{2}.\,\frac{2}{3}.\,\frac{4}{5}.\,60=16$. Die zu $60$ teilerfremden Zahlen sind $1,7,11,13,17,19,23,29,31,37,41,43,47,49,53.\,59.$

Die erste Lösung dieser Aufgabe findet sich in der Abhandlung des berühmten Euler \emph{theoremata arithmetica nova methodo demonstrata}, Comm. nov. Ac. Petrop. VIII, S. 74. Der Beweis wurde später in einer anderen Abhandlung wiederholt: \emph{Speculationes circa quasdam insignes proprietates numerorum}, Acta Petrop. VIII, S. 17.

\Article{39}
Bestimmt man die Bedeutung des Zeichens $\phi$ so, dass $\phi A$ die Anzahl der Zahlen angibt, die zu $A$ teilerfremd und \emph{nicht größer} als $A$ sind, so ist offenbar $\phi1$ nicht mehr $=0$, sondern $=1$. In allen übrigen Fällen ändert sich nichts. Mit dieser Definition erhalten wir den folgenden Satz.

\emph{Sind $a,a',a''$ usw. sämtliche Teiler von $A$, einschließlich der Eins und $A$ selbst, so ist}
\[\phi a+\phi a'+\phi a''+\textit{usw.}=A\]

\emph{Beispiel.} Sei $A=30$, dann ist $\phi1+\phi2+\phi3+\phi5+\phi6+\phi10+\phi15+\phi30=1+1+2+4+2+4+8+8=30$.

\emph{Beweis.} Man multipliziere alle Zahlen, die zu $a$ teilerfremd und nicht größer als $a$ sind, mit $\frac{A}{a}$, ebenso alle zu $a'$ teilerfremden Zahlen mit $\frac{A}{a'}$ usw. Man erhält $\phi a+\phi a'
+\phi a''+$ usw. Zahlen, die alle nicht größer als $A$ sind. Nun gilt:
\SourceLeaf{37}

1) Alle diese Zahlen sind verschieden. Für die Zahlen, die aus \emph{demselben} Teiler von $A$ hervorgehen, ist dies unmittelbar klar. Wären aus verschiedenen Teilern $M,N$ und zu ihnen jeweils teilerfremden Zahlen $\mu,\nu$ gleiche Zahlen hervorgegangen, wäre also $\frac{A}{M}\mu=\frac{A}{N}\nu$, so folgte $\mu N=\nu M$. Wir dürfen $M>N$ annehmen. Da $M$ zu $\mu$ teilerfremd ist und $\mu N$ teilt, müsste es auch $N$ teilen, die größere Zahl also die kleinere. \emph{Q. E. A.}

2) Unter diesen Zahlen finden sich alle Zahlen $1,2,3....A$. Sei eine beliebige Zahl, die $A$ nicht übersteigt, gleich $t$ und der größte gemeinsame Teiler von $A,t$ gleich $\delta$. Dann ist $\frac{A}{\delta}$ ein Teiler von $A$, zu dem $\frac{t}{\delta}$ teilerfremd ist. Die Zahl $t$ findet sich daher offenbar unter denen, die aus dem Teiler $\frac{A}{\delta}$ hervorgehen.

3) Die Anzahl dieser Zahlen ist also $A$, und somit
\[\phi a+\phi a'+\phi a''+\mathrm{usw.}=A.\quad\emph{Q. E. D.}\]

\Article{40}
\emph{Ist der größte gemeinsame Teiler von $A,B,C,D$ usw. gleich $=\mu$, so lassen sich Zahlen $a,b,c,d$ usw. so bestimmen, dass}
\[aA+bB+cC+\textit{usw.}=\mu\]

\emph{Beweis.} Betrachten wir zunächst nur zwei Zahlen $A,B$, deren größter gemeinsamer Teiler $=\lambda$ sei. Dann ist die Kongruenz $Ax\equiv\lambda\ (\mathrm{mod.}\ B)$ lösbar (Artikel 30). Eine Wurzel sei $\equiv\alpha$, und man setze $\frac{\lambda-A\alpha}{B}=\beta$. Dann ist $\alpha A+\beta B=\lambda$, wie verlangt.

Kommt eine dritte Zahl $C$ hinzu, so habe der größte gemeinsame Teiler der Zahlen $\lambda,C$ den Wert $=\lambda'$. Dieser ist zugleich der größte gemeinsame Teiler von $A,B,C$\HistoricalNote{*)}{Denn offenbar teilt $\lambda'$ alle Zahlen $A,B,C$. Wäre es nicht der \emph{größte} gemeinsame Teiler, so wäre dieser größer als $\lambda'$. Da dieser größte Teiler aber $A,B,C$ teilt, müsste er auch $k\alpha A+k\beta B+\gamma C$, also $\lambda'$ teilen, die größere Zahl also die kleinere. \emph{Q. E. A.} --- Noch leichter folgt dies aus Artikel 18.}. Man bestimme Zahlen $k,\gamma$ so, dass $k\lambda+\gamma C=\lambda'$ gilt; dann ist $k\alpha A+k\beta B+\gamma C=\lambda'$.

Kommt eine vierte Zahl $D$ hinzu, so setze man den größten gemeinsamen Teiler von $\lambda',D$, der offenbar zugleich der größte gemeinsame Teiler von $A,B,C,D$ ist, gleich $=\lambda''$ und bilde $k'\lambda'+\delta D=\lambda''$. Dann ist $kk'\alpha A+kk'\beta B+k'\gamma C+\delta D=\lambda''$.

\SourceLeaf{38}
Ebenso kann man bei beliebig vielen weiteren Zahlen vorgehen.

Haben also die Zahlen $A,B,C,D$ usw. keinen gemeinsamen Teiler, so lässt sich offenbar erreichen, dass
\[aA+bB+cC+\mathrm{usw.}=1\]

\Article{41}
\emph{Ist $p$ eine Primzahl und hat man $p$ Gegenstände, von denen beliebig viele gleich sein dürfen, nur nicht alle, so ist die Anzahl ihrer Permutationen durch $p$ teilbar.}

\emph{Beispiel.} Die fünf Gegenstände $A,A,A,B,B$ lassen sich auf zehn verschiedene Arten anordnen.

Dieser Satz lässt sich leicht aus der bekannten Theorie der Permutationen beweisen. Gibt es unter den Gegenständen zunächst $a$ gleiche, nämlich $=A$, dann $b$ gleiche, nämlich $=B$, dann $c$ gleiche, nämlich $=C$ usw., wobei die Zahlen $a,b,c$ usw. auch gleich eins sein dürfen, so ist
\[a+b+c+\mathrm{usw.}=p\]
Die Anzahl der Permutationen ist dann
\[=\frac{1.2.3......p}{1.2.3...a.1.2..b.1.2..c\ \mathrm{usw.}}\]
Offenbar ist der Zähler dieses Bruchs durch den Nenner teilbar, da die Anzahl der Permutationen eine ganze Zahl sein muss. Der Zähler ist durch $p$ teilbar; der Nenner dagegen, der aus Faktoren kleiner als $p$ zusammengesetzt ist, ist nicht durch $p$ teilbar (Artikel 15). Somit ist die Anzahl der Permutationen durch $p$ teilbar (Artikel 19).

Wir hoffen aber, dass manche Leser auch an dem folgenden Beweis Freude haben werden.

Zwei Permutationen heißen \emph{ähnlich}, wenn sich die Reihenfolge ihrer Gegenstände nur darin unterscheidet, dass der Gegenstand, der in der einen an erster Stelle steht, in der anderen einen anderen Platz einnimmt, die übrigen aber in beiden in derselben Reihenfolge aufeinanderfolgen und auf den letzten Gegenstand wieder der erste folgt\HistoricalNote{*)}{Denkt man sich ähnliche Permutationen im Kreis geschrieben, sodass der letzte Gegenstand an den ersten grenzt, so unterscheiden sie sich überhaupt nicht, da keine Stelle als erste oder letzte bezeichnet werden kann.}. So sind in unserem Beispiel die Permutationen $ABAAB$ und $ABABA$ ähnlich, weil die Gegenstände, die in der ersteren an erster, zweiter usw. Stelle stehen, in der letzteren in derselben Reihenfolge an dritter, vierter usw. Stelle stehen.

\SourceLeaf{39}

Da jede Permutation aus $p$ Gegenständen besteht, kann man zu jeder $p-1$ ähnliche finden, indem man den zuvor ersten Gegenstand an die zweite, dritte usw. Stelle verschiebt. Können keine zwei davon identisch sein, so ist offenbar die Gesamtzahl der Permutationen durch $p$ teilbar, denn sie ist $p$-mal so groß wie die Anzahl aller untereinander unähnlichen Permutationen. Nehmen wir also an, zwei Permutationen

\[PQ\ldots TV\ldots YZ;\qquad V\ldots YZPQ\ldots T\]

\noindent die durch Verschieben der Glieder auseinander hervorgehen, seien identisch, also $P=V$ usw. Das Glied $P$, das in der ersteren an erster Stelle steht, stehe in der letzteren an Stelle $n+1$. Dann ist in der letzteren Reihe das Glied an Stelle $n+1$ dem ersten, das an Stelle $n+2$ dem zweiten gleich usw. Somit ist auch das Glied an Stelle $2n+1$ wieder dem ersten gleich, ebenso das an Stelle $3n+1$ usw. Allgemein ist das Glied an Stelle $kn+m$ dem an Stelle $m$ gleich. Übersteigt dabei $kn+m$ die Zahl $p$, so ist entweder die Reihe $V\ldots YZPQ\ldots T$ immer wieder von vorn zu beginnen, oder von $kn+m$ ist das nächstkleinere Vielfache von $p$ abzuziehen. Bestimmt man nun $k$ so, dass $kn\equiv1\ (\mathrm{mod.}p)$ gilt, was möglich ist, da $p$ eine Primzahl ist, so ist allgemein das Glied an Stelle $m$ dem an Stelle $m+1$ gleich. Jedes Glied ist also dem folgenden gleich, \emph{das heißt}, alle Glieder sind gleich, entgegen der Voraussetzung.

\Article{42}

\emph{Sind die Koeffizienten $A,B,C\ldots.N$, $a,b,c\ldots.n$ zweier Funktionen der Form}

\[
\begin{array}{l@{\qquad}c}
x^m+Ax^{m-1}+Bx^{m-2}+Cx^{m-3}\mathinner{\ldotp\ldotp\ldotp\ldotp\ldotp}+N
\quad\mathinner{\ldotp\quad\ldotp\quad\ldotp\quad\ldotp\quad\ldotp}&(P)\\
x^\mu+ax^{\mu-1}+bx^{\mu-2}+cx^{\mu-3}\mathinner{\ldotp\ldotp\ldotp\ldotp\ldotp}+n
\quad\mathinner{\ldotp\quad\ldotp\quad\ldotp\quad\ldotp\quad\ldotp}&(Q)
\end{array}
\]

\noindent\emph{sämtlich rational, aber nicht alle ganzzahlig, und ist das Produkt von $(P)$ und $(Q)$ gleich}

\[=x^{m+\mu}+\mathfrak{A}x^{m+\mu-1}+\mathfrak{B}x^{m+\mu-2}+\emph{usw.}+\mathfrak{Z}\]

\noindent\emph{so können nicht alle seine Koeffizienten $\mathfrak{A},\mathfrak{B}\ldots.\mathfrak{Z}$ ganzzahlig sein.}

\emph{Beweis.} Man schreibe alle Brüche unter den Koeffizienten $A,B$ usw., $a,b$ usw. vollständig gekürzt und wähle eine beliebige Primzahl $p$, die einen oder mehrere ihrer Nenner teilt. Wir dürfen annehmen, dass $p$ den Nenner eines gebrochenen Koeffizienten in $(P)$ teilt. Wird $(Q)$ durch $p$ geteilt, so gibt es auch in $\frac{(Q)}{p}$ mindestens einen gebrochenen Koeffizienten, dessen Nenner den Faktor $p$ enthält, nämlich den ersten Koeffizienten $\frac{1}{p}$.

\SourceLeaf{40}

In $(P)$ gibt es offenbar einen gebrochenen Term, dessen Nenner eine \emph{höhere} Potenz von $p$ enthält als die Nenner aller vorhergehenden gebrochenen Terme und eine \emph{nicht niedrigere} als die Nenner aller folgenden. Dieser Term sei $=Gx^g$, und die Anzahl der Faktoren $p$ im Nenner von $G$ sei $=t$. Einen entsprechenden Term gibt es in $\frac{(Q)}{p}$; er sei $=\varGamma x^\gamma$ und die Anzahl der Faktoren $p$ im Nenner von $\varGamma$ sei $=\tau$. Offenbar ist $t+\tau$ mindestens $=2$. Nach diesen Vorbereitungen hat der Term $x^{g+\gamma}$ im Produkt von $(P)$ und $(Q)$ einen gebrochenen Koeffizienten, dessen Nenner $t+\tau-1$ Faktoren $p$ enthält. Dies wird folgendermaßen bewiesen.

Die Terme, die in $(P)$ dem Term $Gx^g$ vorausgehen, seien ${}'Gx^{g+1}$, ${}''Gx^{g+2}$ usw., die folgenden dagegen $G'x^{g-1}$, $G''x^{g-2}$ usw. Ebenso seien in $\frac{(Q)}{p}$ die dem Term $\varGamma x^\gamma$ vorausgehenden Terme ${}'\varGamma x^{\gamma+1}$, ${}''\varGamma x^{\gamma+2}$ usw., die folgenden $\varGamma'x^{\gamma-1}$, $\varGamma''x^{\gamma-2}$ usw. Dann ist im Produkt von $(P)$, $\frac{(Q)}{p}$ der Koeffizient von $x^{g+\gamma}$ gleich

\[
\begin{aligned}
&=G\varGamma+{}'G\varGamma'+{}''G\varGamma''+\text{usw.}\\
&\phantom{=G\varGamma}\mathord{+}{}'\varGamma G'+{}''\varGamma G''+\text{usw.}
\end{aligned}
\]

Das Glied $G\varGamma$ ist ein Bruch, dessen Nenner bei vollständiger Kürzung $t+\tau$ Faktoren $p$ enthält. Die übrigen Glieder enthalten, soweit sie gebrochen sind, weniger Faktoren $p$ im Nenner. Denn sie sind Produkte zweier Faktoren, von denen der eine nicht mehr als $t$, der andere weniger als $\tau$ Faktoren $p$ enthält, oder der eine nicht mehr als $\tau$ und der andere weniger als $t$. Somit hat $G\varGamma$ die Form $\frac{e}{fp^{t+\tau}}$, die Summe der übrigen Glieder dagegen die Form $\frac{e'}{f'p^{t+\tau-\delta}}$, wobei $\delta$ positiv ist und $e,f,f'$ frei vom Faktor $p$ sind. Die Summe aller Glieder ist daher $=\frac{ef'+e'fp^\delta}{ff'p^{t+\tau}}$. Ihr Zähler ist nicht durch $p$ teilbar, sodass der Nenner durch keine Kürzung weniger als $t+\tau$ solche Faktoren enthalten kann. Der Koeffizient von $x^{g+\gamma}$ im Produkt von $(P),(Q)$ ist also

\[=\frac{ef'+e'fp^\delta}{ff'p^{t+\tau-1}}.\]

\noindent\emph{das heißt, ein Bruch}, dessen Nenner $t+\tau-1$ Faktoren $p$ enthält.\par
\noindent\mbox{\emph{Q. E. D.}}

\SourceLeaf{41}

\Article{43}

\emph{Eine Kongruenz vom Grad $m$}

\[Ax^m+Bx^{m-1}+Cx^{m-2}+\emph{usw.}+Mx+N\equiv0\]

\noindent\emph{deren Modul eine Primzahl $p$ ist, die $A$ nicht teilt, lässt sich nicht auf mehr als $m$ verschiedene Arten lösen; sie hat also nicht mehr als $m$ nach $p$ inkongruente Wurzeln} (siehe Artikel 25, 26).

Wer dies bestreitet, nehme an, es gebe Kongruenzen verschiedener Grade $m,n$ usw., die mehr als $m,n$ usw. Wurzeln haben. Der kleinste solche Grad sei $m$, sodass alle entsprechenden Kongruenzen niedrigeren Grades unserem Satz genügen. Da dies für den ersten Grad bereits bewiesen ist (Artikel 26), muss $m$ entweder $=2$ oder größer sein. Die Kongruenz

\[Ax^m+Bx^{m-1}+\text{usw.}+Mx+N\equiv0\]

\noindent hat also mindestens $m+1$ Wurzeln, etwa $x\equiv\alpha,x\equiv\beta,x\equiv\gamma$ usw. Wir dürfen annehmen, dass alle Zahlen $\alpha,\beta,\gamma$ usw. positiv und kleiner als $p$ sind und $\alpha$ die kleinste ist. Setzt man nun in der gegebenen Kongruenz für $x$ den Ausdruck $y+\alpha$ ein, so geht sie über in

\[A'y^m+B'y^{m-1}+C'y^{m-2}....+M'y+N'\equiv0\]

Diese Kongruenz ist offenbar erfüllt, wenn man $y\equiv0$ oder $\equiv\beta-\alpha$ oder $\equiv\gamma-\alpha$ usw. setzt. Diese Wurzeln sind alle verschieden, und ihre Anzahl ist $=m+1$. Da $y\equiv0$ eine Wurzel ist, folgt, dass $N'$ durch $p$ teilbar ist. Also gilt auch

\[y(A'y^{m-1}+B'y^{m-2}+\text{usw.}+M')\text{ ist }\equiv0\ (\mathrm{mod.}p),\]

\noindent wenn $y$ einer der $m$ Werte $\beta-\alpha,\gamma-\alpha$ usw. gegeben wird, die alle $>0$ und $<p$ sind. In all diesen Fällen gilt daher auch

\[A'y^{m-1}+B'y^{m-2}+\text{usw.}+M'\text{ ist }\equiv0\ \text{(Artikel 22)}\]

\noindent\emph{das heißt}, die Kongruenz \qquad $A'y^{m-1}+B'y^{m-2}+\text{usw.}+M'\equiv0$

\noindent hat den Grad $m-1$ und $m$ Wurzeln und widerspricht somit unserem Satz. Denn offenbar ist $A'$ gleich $=A$ und daher, wie verlangt, nicht durch $p$ teilbar. Dabei hatten wir angenommen, dass alle Kongruenzen von niedrigerem Grad als $m$ dem Satz genügen. \mbox{\emph{Q. E. A.}}

\SourceLeaf{42}

\Article{44}

Obwohl wir hier vorausgesetzt haben, dass der Modul $p$ den Koeffizienten des höchsten Terms nicht teilt, ist der Satz nicht auf diesen Fall beschränkt. Wäre nämlich der erste Koeffizient oder wären auch einige der folgenden durch $p$ teilbar, so könnten diese Terme ohne Weiteres wegfallen. Dadurch würde die Kongruenz schließlich auf einen niedrigeren Grad zurückgeführt, bei dem der erste Koeffizient nicht mehr durch $p$ teilbar wäre. Dies gilt, sofern nicht alle Koeffizienten durch $p$ teilbar sind; dann wäre die Kongruenz eine Identität und die Unbekannte völlig unbestimmt.

Dieser Satz wurde zuerst vom berühmten La Grange aufgestellt und bewiesen (\emph{M\'em. de l'Ac. de Berlin, Ann\'ee} 1768, \emph{S.} 192). Er findet sich auch in der Abhandlung des berühmten Le Gendre \emph{Recherches d'Analyse ind\'etermin\'ee, Hist. de l'Acad. de Paris} 1785, \emph{S.} 466. Der berühmte Euler bewies in \emph{Nov. Comm. Ac. Petr.} XVIII, \emph{S.} 93, dass die Kongruenz $x^n-1\equiv0$ nicht mehr als $n$ verschiedene Wurzeln haben kann. Obwohl dies ein Sonderfall ist, lässt sich die Methode, die dieser große Mann verwendet hat, leicht auf alle Kongruenzen übertragen. Einen noch eingeschränkteren Fall hatte er schon zuvor behandelt, in \emph{Comm. nov. Ac. Petr.} V, \emph{S.} 6; diese Methode lässt sich aber nicht allgemein anwenden. Weiter unten, in Abschnitt VIII, werden wir den Satz noch auf eine andere Weise beweisen. So verschieden alle diese Methoden auf den ersten Blick erscheinen mögen, werden kundige Leser bei einem Vergleich leicht erkennen, dass sie alle auf demselben Grundsatz beruhen. Da der Satz hier nur als Hilfssatz betrachtet wird und eine vollständige Darstellung nicht hierhergehört, verzichten wir auf eine gesonderte Behandlung zusammengesetzter Moduln.

\BookRule{0.28\linewidth}
